% Frontier-closure search — describe the two-phase algorithm, the
% in-memory overlay, the dependency-ordered materialise pass, and
% the verification pass. This is the algorithmic contribution of
% the paper.

The frontier-closure search converges the catalog to a fixed point
under the derivation operators of Section \ref{sec:strategies}. It
propagates rank improvements along a dependency DAG of shapes
(below) rather than re-deriving everything from scratch; the loop
terminates when no shape improves.

\subsection{Two-phase structure}

A naive closure interleaves \emph{search} (find the best
recombination) and \emph{materialise} (build and write the concrete
\((U, V, W)\)) at every shape, every round. We split them:

\begin{description}[itemsep=2pt]
  \item[Phase 1 -- Search.] Iterate rounds. For each shape, run the
    derivation search to find the best predicted rank
    \(\hat r = \sum_k r_{\text{sub-}k}\). Strict improvements over the
    shape's current best (catalog \emph{or} overlay) are recorded in
    an in-memory \emph{pending overlay}, keyed by shape. The overlay
    is consulted by the rank resolver in subsequent rounds, so a win
    at round \(N\) propagates to round \(N+1\) without writing
    anything to disk. Repeat until a round produces no new overlay
    entries.
  \item[Phase 2 -- Materialise.] Iterate the overlay in
    dependency order (smaller maximum-dimension first). For each
    entry, materialise the strategy via
    \code{Recombination.constructWithAllocation} or its concat /
    Kronecker analogue. Sub-product lookups are served from a
    composed \code{AlgorithmLookup} that returns
    just-materialised winners alongside the on-disk catalog, so a
    larger winner can use a smaller winner from this same phase as
    a sub-product. Spot-check the result against the predicted
    rank, verify, write to disk.
\end{description}

Theorem \ref{thm:nonoverlap} guarantees that the rank propagated by
Phase 1 equals the rank produced by Phase 2 for every entry that
materialises successfully (the rare failures are infrastructural --
e.g.\ a sub-product that became unavailable between rounds).

\paragraph{The impact DAG, not a full sweep.} Shapes and their
derivation candidates form a directed acyclic graph of \emph{impacts}:
an edge \(\sigma \to \tau\) exists whenever a rank improvement at shape
\(\sigma\) can lower the predicted rank at \(\tau\) --- because \(\tau\)
Kron-, concat- or recombines a copy of \(\sigma\) (an \emph{upward} edge),
or because \(\tau\) is a projection of \(\sigma\) (a \emph{downward}
edge; Section~\ref{ssec:recombination}). A win therefore cannot change
any shape that is not downstream of it. The search exploits exactly this:
the frontier driver (\code{FrontierClosure}) keeps a priority queue of
seeds keyed by \(\omega\) (best first), pops a seed, and on each
improvement pushes back \emph{only} the shapes that seed can impact ---
so it rides the DAG, touching a shape only when one of its sub-shapes has
just improved. We do \emph{not} perform a full Cartesian sweep of all
shapes against all strategies every round: work is proportional to the
impactful edges actually traversed, not to
\(|\text{shapes}| \times |\text{strategies}| \times |\text{rounds}|\).
The two cost tiers of this DAG also differ --- an upward edge is a cheap
rank-arithmetic re-evaluation, whereas a downward (projection) edge
requires materialising and DCE-ing the parent --- so the closure fans out
upward eagerly and schedules the downward work.

\subsection{Why the split matters}

Three reasons.

\paragraph{Throughput.} Materialisation involves Java array
allocation, factor-matrix arithmetic, spot-check sampling, and disk
I/O. A target shape at \(\nmpshape{17}{17}{17}\) takes a fraction of
a second to predict and a few seconds to materialise; at
\(\nmpshape{32}{32}{32}\) the spread becomes minutes. Postponing
materialisation lets the search round terminate as soon as no new
predictions are found, even when 50 shapes have pending materialise
work.

\paragraph{Correctness of intra-round propagation.}
A previous incarnation of the closure loop wrote each materialised
winner to disk and relied on the next round to reload it. The
catalog \code{FieldAwareLookup} was loaded once at the start of the
sweep, so the next round did \emph{not} see the new entry: round
\(N+1\) used the same catalog snapshot as round \(N\). The pending
overlay fixes this: subsequent rounds see prior-round wins
immediately, regardless of whether the JVM was restarted.

\paragraph{Decoupling verification.} Verification is a third phase
that can run on demand against any materialised scheme -- e.g.\ as a
periodic catalog-wide sweep or as a follow-up to a search batch.
There is no need to verify every newly-materialised scheme inside
the closure loop, only those whose predicted-vs-actual rank match
proves the materialise step succeeded (which is the cheap check;
true random-matmul spot-check can run in batch later).

\subsection{Rebuilding the catalog from atoms: why the closure is multipass}
\label{ssec:multipass}

It helps to state the closure's job in the strongest form: given only the
\emph{atoms} (Section \ref{sec:architecture}) --- the handful of primitive
schemes we do not derive --- reconstruct the entire catalog by repeatedly
applying the derivation operators of Section \ref{sec:strategies}. The
round-based fixed point is exactly what makes this safe and easy: rather than
hand-ordering which shape depends on which, each round re-applies every
operator to every shape and the loop stops only when nothing improves.

For an \emph{upward}-only operator set this is almost an accounting detail.
Kronecker, axis concatenation, recombination and the serendipitous product
(Section \ref{ssec:serendipitous}) all build a \emph{larger} shape from
\emph{smaller} ones, so a single sweep in increasing maximum dimension would
already suffice: every sub-product a shape needs is final before the shape is
processed.

The multipass design earns its keep once an operator runs \emph{downward}.
The canonical example is \emph{projection}: a scheme for
\(\nmpshape{n}{m}{p}\) is obtained from a \emph{larger} one for, say,
\(\nmpshape{n{+}1}{m}{p}\) by zeroing one input slice, discarding the matching
output slice, and dead-code-eliminating the products that fed only the dropped
index. All three axes project. The two outer axes ($n$, $p$) each touch one
input operand and the output $C$, so projecting them shrinks $C$ (a row or a
column). The inner/contracted axis $m$ is different: it indexes A's columns
\emph{and} B's rows but \emph{not} $C$, so projecting $m$ removes a term from
the contraction sum --- $C$ keeps its shape, and the dead-code test consults
\emph{both} input factors rather than one input and the output. A cube-to-cube
step such as \(\nmpshape{26}{26}{26}\to\nmpshape{25}{25}{25}\) projects one
index on each of the three axes at once (the index choice per axis is what
FMM's \([[1,15],[15]]\) bracket records). Two facts make this concrete:
\begin{itemize}[itemsep=2pt]
  \item \(\nmpcount{}{3}{3}{3}{25}\) is the projection of
        \(\nmpcount{}{4}{4}{4}{49}\): pad \(\nmpshape{3}{3}{3}\) with a zero
        row and column to \(\nmpshape{4}{4}{4}\), apply Strassen recursively,
        and eliminate the padding --- precisely the worked example of
        \cite{drisc09}. Projection is thus the same operation as
        padding-plus-DCE, run in reverse.
  \item At scale, \(\nmpcount{}{25}{25}{25}{8359}\) is the projection of
        \(\nmpcount{}{26}{26}{26}{8658}\); FMM \cite{fmmlille} builds most of
        its non-cube and odd-cube entries this way, projecting a strong even
        cube onto its neighbours.
\end{itemize}

Which index is dropped matters: each index has a different set of products
localised to it, so dropping the most-localised one eliminates the most and
gives the lowest rank. A concrete case --- projecting Laderman's
\(\nmpcount{}{3}{3}{3}{23}\) down one axis to \(\nmpshape{2}{3}{3}\):
\begin{center}
\begin{tabular}{lcc}
\toprule
dropped index & products eliminated (DCE) & resulting rank \\
\midrule
\(0\) & 5 & 18 \\
\(1\) & 7 & \(16\) \\
\(2\) & 7 & \(16\) \\
\bottomrule
\end{tabular}
\end{center}
Same scheme, same target shape, but a two-multiplication swing from the choice
alone: Laderman concentrates more products on indices 1 and 2 than on index 0,
so dropping either of those eliminates 7 products versus only 5 for index 0. A
``drop-the-last-index'' rule would have left rank 18 where 16 was available.
(The format's true optimum, \(R(\nmpshape{2}{3}{3})=15\), comes from a
different scheme, not a projection of Laderman --- a reminder that projection's
payoff is bounded by its parent, which is why the closure keeps several parents
per format.)

FMM records this choice explicitly (e.g.\ the bracket in
``projection \([[1,15],[15]]\)'' selects index 15, not the last index), and
the operator must try every index per axis rather than only the tail. The best
choice is moreover joint in the source scheme and the dropped index --- a
different parent for the same shape concentrates products differently --- so
projection, like the serendipitous product, rewards keeping structurally
diverse schemes per format rather than only the minimum-rank one. With our
recursively-composed schemes the per-index localisation is not predictable in
closed form, so we simply enumerate all drop positions, project, and verify ---
each position is an \(O(r)\) support scan plus dead-code elimination, so the
exhaustive choice is cheap.

A downward operator destroys the acyclicity in dimension: improving
\(\nmpshape{26}{26}{26}\) can lower \(\nmpshape{25}{25}{25}\), which a
smaller-first sweep has already finalised. So the catalog cannot be rebuilt in
one dimension-ordered pass --- it must iterate to a fixed point, re-projecting
neighbours after each cube improvement. This is also the \emph{safer} design:
instead of special-casing ``did some larger shape just improve?'', the
fixed point simply re-runs all operators each round and a projection win is
absorbed on the next round automatically. Cube quality therefore propagates
downward through the whole catalog --- a single better \(\nmpshape{n}{n}{n}\)
improves every nearby rectangular and odd-cube shape at once.

The downward (projection) operator and its \emph{projection-margin} score
\(\mu\) --- the criterion that decides when a higher-rank parent projects to
a lower-rank child --- are defined alongside the derivation strategies in
Section~\ref{ssec:projmargin}. Operationally the closure fires a bounded
\emph{downward wave} on each catalog update: a freshly arrived scheme with
\(\mu>0\) is projected into its dimension-neighbours, and any neighbour that
improves is enqueued for its own projections; termination is immediate since
dimension strictly decreases.

\subsection{Pool, allocation enumeration, and pruning}

The derivation search at a shape \(\nmpshape{n}{m}{p}\) enumerates,
over a pool of outer schemes:
\begin{itemize}[itemsep=2pt]
  \item per-axis allocations of \(n / m / p\) into the outer scheme's
    block counts, optionally filtered by a maximum per-allocation
    imbalance \(\max(\mathbf{a}) - \min(\mathbf{a}) \le
    \texttt{maxImbalance}\) and capped at
    \texttt{maxCombinations} via a top-\(K\)-by-balance pre-selector;
  \item over-allocation by per-axis padding \(\delta\) up to
    \texttt{maxPadding}, paired with a corresponding tail-peel
    pattern (Section \ref{ssec:peel}); and
  \item axis-flip orbit masks of the outer scheme, evaluated
    analytically via per-product block-support bitmasks
    (\code{AnalyticalMaskSearch}) so the 8-mask coverage costs only
    a constant factor over single-mask evaluation.
\end{itemize}

The default pool (\code{PoolConfig.simple}) carries a small number of
high-value outer schemes -- Strassen, Winograd, Smirnov 3-row
templates, and a few imported Pan / AT / FMM bases. An
\code{extendedPool} mode adds every \code{Leaf}-tagged catalog
entry as an outer template; that pool is large enough that the
maxCombinations-capped allocation enumerator becomes load-bearing.

\subsection{Allocation optimisation by branch-and-bound}\label{ssec:alloc-bnb}

Enumerating allocations naively is hopeless: for a fixed outer scheme of
shape \(\nmpshape{n_o}{m_o}{p_o}\), the number of ways to split the target
is \(\binom{N-1}{n_o-1}\binom{M-1}{m_o-1}\binom{P-1}{p_o-1}\), which blows
up well before \(N = 32\). Yet the inner problem --- given a target,
template, and orientation, find the allocation minimising
\(\sum_{\sigma} R(\sigma)\) over the sub-problem multiset (Section
\ref{ssec:decomposition}) --- is solved \emph{exactly} by
\code{AllocationOptimizer}, a branch-and-bound over the allocation grid.

The search fixes the three axis partitions in turn (\(\mathbf{a}_N\),
then \(\mathbf{a}_M\), then \(\mathbf{a}_P\)), depth-first. Two ingredients
make it cheap:
\begin{itemize}[itemsep=2pt]
  \item \textbf{A strong incumbent.} The Kronecker (uniform-allocation)
    rank is a ready upper bound, seeded before the descent; most subtrees
    then fall to the bound immediately.
  \item \textbf{An admissible lower bound.} At a partial node --- some
    axes resolved, others open --- each product is charged the smallest
    rank it could still take (the resolved axes' sizes, a floor on the
    open ones). Because \(R\) is monotone in the dimensions this never
    overestimates, so a node whose bound already meets the incumbent is
    pruned without expansion.
\end{itemize}
Branch-and-bound returns the \emph{same} optimum as the exhaustive grid
while visiting far fewer allocations (the result reports
\code{nodes} visited against \code{fullSpace}).

\paragraph{Budget and what ``optimal'' means here.} The optimiser runs
under a \code{SearchBudget}: \code{EXACT} (no cap), \code{upTo}(a known
upper bound, e.g.\ the Kronecker rank), or a \code{maxNodes} anytime cap.
The result carries an \code{exhaustive} flag, and we propagate it
honestly (cf.\ the optimality discipline of Section~\ref{sec:intro}). When
\code{exhaustive} is true the returned allocation is the \emph{proven
global minimum over allocations} for that template and orientation ---
\emph{optimal-within-scope}, not the tensor rank \(R(\nmpshape{N}{M}{P})\),
which a different template or a non-recombination construction may still
beat. When the node budget is hit, \code{exhaustive} is false and the
rank is only a best-found \emph{upper bound}. The outer loop (pool of
templates \(\times\) axis-flip orbit) and the frontier closure wrap this
exact inner solve, and are themselves bounded; the catalog never reports
an allocation result as a global rank optimum.

\subsection{Sound pruning from cheap candidates}\label{ssec:cheap-pruning}

The derivation search has three families of candidate strategies
at every target shape: multi-base recombination (heavy --- iterates
the entire pool, every allocation, every axis-flip mask),
\emph{concat-split} (cheap --- a single axis split, two SOTA
lookups), and \emph{Kronecker} (cheap --- nested factor enumeration
of \(n, m, p\), all SOTA lookups). Cheap candidates finish in
milliseconds; the heavy recombination sweep can take minutes on a
single shape at high dim.

We exploit this asymmetry. Concat and Kronecker run \emph{first},
producing an upper bound \(\hat r_{\text{cheap}} = \min(r_{\text{concat}},
r_{\text{kron}})\). The recombination sweep is then seeded with
\(\text{bestRank} := \hat r_{\text{cheap}}\). Two sound prunes follow:

\begin{itemize}[itemsep=2pt]
  \item \emph{Per-base lower bound.} For an outer base of rank \(r\),
    the recombination produces \(r\) sub-products, each costing
    \(\ge 1\); the pair-fusion pass (Section \ref{ssec:pair-fusion})
    only applies pairs whose savings are positive, so pair-cost is
    at least the unpaired sum on degenerate-shape sub-products.
    Hence \(\text{total} \ge r\) under any allocation. If
    \(r \ge \text{bestRank}\), the entire base is skipped.
  \item \emph{Best-so-far gate.} After each
    \code{(allocation, mask)} tuple completes, the predicted rank is
    compared against the running \(\text{bestRank}\), updating the
    candidate only on strict improvement.
\end{itemize}

We deliberately do \emph{not} prune \emph{inside} the sub-product
accumulation loop. A running sum above \(\text{bestRank}\) does not
imply the paired total is above it, since pair-fusion can reduce
the sum by tens of percent on cyclic-shape multisets. Adding a
slack constant would be a heuristic, and we keep the search at
SAT-style soundness: every prune is provably correct.

A canonical case where this matters: at \(\nmpshape{12}{12}{12}\),
the Kronecker pair \(\nmpshape{2}{4}{4}{=}26 \otimes
\nmpshape{6}{3}{3}{=}40\) gives 1040 in milliseconds. Without the
reorder, the recombination sweep grinds for several minutes on the
AlphaEvolve \(\nmpshape{5}{5}{5}{=}93\) base, finding nothing
better than 1071, before the final \(\min\) over all strategies
selects Kronecker anyway. With the reorder, AE-\(\nmpshape{5}{5}{5}\)
still runs (its trivial lower bound \(r=93 < 1040\)), but every
allocation that doesn't beat 1040 is silently dropped instead of
being recorded as a (worthless) intermediate best.

\subsection{Low-\(\omega\) contamination}\label{ssec:omega-contamination}

A complementary optimisation reshapes the \emph{order} in which
shapes are processed during a closure round. The naive iteration is
lexicographic over \((n, m, p)\), so the search reaches
\(\nmpshape{12}{12}{12}\) before it reaches \(\nmpshape{17}{17}{17}\)
and so on. But the value of an early win depends on its implied
\(\omega\): a low-\(\omega\) seed contaminates every larger shape
that consumes it as a sub-product.

We rank candidate outer bases (and known catalog shapes) by their
implied exponent
\(\omega(\nmpshape{n}{m}{p}{:}r) = \log_n(r) \cdot 3 / (1 + \log_n(m)
+ \log_n(p))\)
under the convention of Section \ref{sec:notation}, and prioritise:

\begin{enumerate}[itemsep=2pt]
  \item Low-\(\omega\) seeds first: e.g.\
    \(\nmpshape{4}{4}{4}{=}48\) (\(\omega \approx 2.7925\); AlphaEvolve
    over \(\Complex\), the DPS rationalisation \cite{dps2025} over
    \(\Rationals/\Reals\)),
    Smirnov \(\nmpshape{3}{3}{6}{=}40\), Pan
    \(\nmpshape{3}{4}{7}{=}63\), then up the ladder.
  \item Targets that can be \emph{contaminated} by a low-\(\omega\)
    seed (via Kronecker, recombination, or concat) are re-evaluated
    in the same round once the seed is on the overlay, so a
    fresh \(\omega\) reduction propagates without a separate round.
  \item The shape order within each round respects this dependency:
    every shape that consumes a seed via a known composition rule
    waits for the seed's seed-round completion before its own
    search starts.
\end{enumerate}

Concretely: when AE-\(\nmpshape{4}{4}{4}{=}48\) lands as a seed in
round \(k\), every Kronecker target whose factorisation includes
\((4, 4, 4)\) on at least one axis pair --- \(\nmpshape{8}{8}{8}\),
\(\nmpshape{12}{12}{12}\), \(\nmpshape{16}{16}{16}\), and so on ---
sees its predicted rank drop in round \(k+1\). The ``contamination''
metaphor captures the spreading: a low-\(\omega\) reduction at
\(\nmpshape{4}{4}{4}\) silently rewrites everyone's upper bound.

The selection of which schemes act as seeds is encoded in
\code{docs/bases-by-omega.md}, an automatically-regenerated table
ranking every pool-eligible base by its raw \(\omega\), its
\(\omega\) after imbalance-\(\le 3\) restriction, and its \(\omega\)
after axis-flip-canonical and S\(_3\)-canonical orbit reductions.
The lowest-\(\omega\) row at the reference target
\(\nmpshape{17}{17}{17}\) is the natural seed for any sweep that
targets dim-32-or-higher cubic shapes.

\subsection{Pruning is opt-in}

A reviewer's natural question is whether the search uses heuristics
to bound the allocation space. We have a strong project-level
convention that pruning is opt-in: any heuristic that drops part of
the allocation space (\texttt{balancedOnly}, top-\(N\), ``ignore
imbalanced''---) is a flag the user explicitly turns on, not a
silent default. This is a reaction to a recurring failure mode where
``the search did not find scheme X'' turns out to mean ``the search
silently pruned the allocation that would have found scheme X.''
Unbalanced cubic targets and bridge-shape sub-problems (Hopcroft--Kerr
type) live precisely in the region that naive balanced-only pruning
discards.

\subsection{Caching}

Caching exists at two layers. The first is per-scheme \emph{symbolic
templates}: for each outer scheme we extract once a tuple of six
block-support bitmasks per product (\(U\)-row, \(U\)-col, \(V\)-row,
\(V\)-col, \(W\)-row, \(W\)-col), and per-call sub-shape extraction
reduces from \(O(r \cdot \text{dim}^2)\) to
\(O(r \cdot \text{dim})\). The second is a per-(scheme, allocation,
peel) cache of the resulting sub-shape multiset; this is a
hashing-overhead trade and is currently capped at 200k entries per
scheme. A symbolic-template-at-allocation-pattern layer would
generalise the second cache; the current implementation captures the
common cases without paying the templating cost.

\subsection{Progress reporting}

A frontier closure over the entire catalog up to dim 32 takes hours
of wall-clock. Per CLAUDE.md ``Long-running procedures'' we emit
periodic progress lines at three nested levels: per-round (round
counter + overlay size), per-shape (search/materialise phase
markers with timing), and intra-search (current pool entry, current
allocation, allocations explored, throughput, best so far). The
intra-search level is the one that catches a stuck inner-loop call,
which otherwise stays silent for tens of minutes during heavy axis-flip
mask enumeration at high dim.
